\documentclass[11pt]{article}
\usepackage[utf8]{inputenc}
\usepackage[T1]{fontenc}
\usepackage{lmodern}
\usepackage[margin=1in]{geometry}
\usepackage{amsmath,amssymb,amsthm,mathtools}
\usepackage{booktabs}
\usepackage{enumitem}
\usepackage[colorlinks=true,linkcolor=black,citecolor=black,urlcolor=black]{hyperref}
\usepackage{longtable}
\usepackage{array}
\usepackage{graphicx}
\usepackage{microtype}
\setlength{\parskip}{0.65em}
\setlength{\parindent}{0pt}

\title{Paper 6 --- REPRESENTATIONS (Dynamics \& Geometry): Dynamics and Geometry as Optimal Admissible Reallocation}
\author{E. Brooke}
\date{}

\begin{document}
\maketitle
\begin{abstract}
This paper derives dynamics and geometry as representational structures that emerge when
 admissibility constraints permit smooth, locally additive bookkeeping. Equations of motion arise
 only in restricted regimes; outside these regimes---at saturation, horizons, or
 measurement---variational dynamics ceases to be meaningful, though admissibility accounting
 remains well-defined. Geometry emerges by a parallel route: distance is minimal correlation cost,
 curvature arises from capacity gradients, and Einstein’s equations appear as the unique closure
 relation in four dimensions (Lovelock uniqueness).

 This February 2026 release incorporates the complete $\Delta$_geo closure: all five geometric
 prerequisites (A9.1--A9.5) are now derived from the axioms, with the metric tensor proved from
 non-factorization (T7B), d = 4 spacetime dimensions derived from capacity budget (T8),
 Lorentzian signature from A4 irreversibility, and the cosmological constant identified as global
 capacity residual (T11). The gravity sector is no longer hypothetical---it is structurally closed.
\end{abstract}

\section*{Part I: Dynamics}



\section{Introduction}



\subsection{The Standard View}

In conventional physics, dynamics and geometry are taken as primitive. Equations of motion are

postulated; spacetime is assumed as an arena. The questions “why do physical laws take

variational form?” or “why is gravity geometric?” are rarely asked.



\subsection{The Constraint-First Inversion}

This paper asks: Under what conditions do equations of motion exist at all? Under what conditions

does geometric description become meaningful? What happens when those conditions fail? The

central claim is that dynamics and geometry are regime-dependent representations of admissible

correlation reallocation. They emerge when correlation capacity admits smooth redistribution,

enforcement cost is locally additive, and saturation boundaries are not encountered.



\section{Why Admissibility Does Not Determine Evolution}

Admissibility determines what correlation structures are jointly enforceable. It specifies necessary

conditions for physical meaningfulness. But it does not provide sufficient conditions for selection.

Given an admissible state, there may be multiple admissible continuations. Admissibility alone

does not pick one.



\section{Regime R: When Dynamics Becomes Meaningful}

Dynamics becomes meaningful when: (R1) Enforcement cost is a smooth function of correlation

sets. (R2) Cost is locally additive over interfaces. (R3) Admissible paths form a connected space.

(R4) No saturation is encountered along the path. When all four conditions hold, enforcement cost

defines a smooth functional on the space of admissible paths, and variational principles emerge

necessarily.



\section{Extremal Paths and Euler-Lagrange Structure}

\[In Regime R, the accumulated enforcement cost along a path defines an action functional: K = $\\int$ L\]

dt, where L encodes the local enforcement cost density. Extremal paths---those minimizing total

\[enforcement cost---satisfy the Euler-Lagrange equations: d/dt($\\partial$L/$\\partial$q) - $\\partial$L/$\\partial$q = 0. This is not\]

postulated. It follows from the variational structure of smooth enforcement cost.



Conservation laws follow by Noether’s theorem: symmetries of the cost functional generate

conserved quantities. Energy conservation follows from time-translation invariance of enforcement

cost. Momentum conservation follows from spatial translation invariance.



\section{Where Dynamics Fails}

Mode                  Cause                              Result



\[Saturation            S_$\Gamma$ = C_$\Gamma$                          Loss of admissible variation\]



Branching             Non-unique continuations           No single equation of motion



Record locking        Irreversible commitment            Regime exit to constraint-dominated



Horizons              Persistent saturation              Dynamical inaccessibility



These are not failures of physical law. They mark the limits of variational encoding. The underlying

admissibility structure continues to hold; only the dynamical representation exits its domain.



\section*{Part II: Geometry}



\section{Correlation Cost and Distance}



\subsection{Definition of Distance}

\textbf{Definition 6.1 (Correlation Distance). The distance between regions A and B is the minimal}

irreducible correlation cost required to jointly enforce coordination between them: d(A, B) $\\equiv$

inf_{$\Gamma$: A$\to$B} $\\int$ C d$\Gamma$, where C is the correlation cost density along path $\Gamma$.



Correlation distance satisfies non-negativity, identity of indiscernibles, and the triangle inequality.

These are forced by enforcement cost structure, not assumed.



\subsection{Geodesics as Least-Cost Paths}

Paths minimizing correlation cost define geodesics. Along such paths, relational structure is

maintained with minimal capacity expenditure. Free motion is least-cost coordination. No force law

is required.



\section{The Metric Tensor from Non-Factorization (T7B)}

The February 2026 release closes the metric derivation completely.



\subsection{Non-Factorization Implies Geometric Response}

\textbf{Theorem T7B [P_structural]. At shared (non-factorizing) interfaces, the external feasibility}

functional is quadratic in displacement. By the polarization identity, a functional that is: (a)

non-degenerate (A1: finite capacity implies no zero-cost directions), (b) continuous (Lipschitz

from the finite-bounded cost condition), and (c) defined on tangent vectors at each point

(from T_continuum smooth structure) is a metric tensor g_$\mu$$\nu$.



This is a definition-level result: a non-degenerate symmetric bilinear form on the tangent space at

each point IS a metric tensor. The framework does not need to postulate geometry---it computes it

from quadratic cost at shared interfaces.



\subsection{Lorentzian Signature (-,+,+,+)}

A4 (irreversibility) implies a strict partial order on admissible configurations. By the

Hawking-King-McCarthy theorem (1976), a strict partial order on events satisfying appropriate

density conditions determines the conformal structure. Combined with Malament’s theorem (1977)

\[and volume normalization ($\Omega$ = 1 by Radon-Nikodym uniqueness), this fixes the Lorentzian\]

signature. The metric has exactly one timelike direction. [Imports: HKM 1976, Malament 1977]



\section{Spacetime Dimension d = 4 (T8)}

\textbf{Theorem T8 [P_structural]. d = 4 is the minimal admissible spacetime dimension. d $\leq$ 3 is}

\[hard-excluded: d = 3 has zero propagating graviton degrees of freedom (no gravitational\]

records, violating A4); d = 2, 3 lack knots for topological stability of bound states. d = 4:

minimum dimension with propagating gravity (2 DOF) + knot stability + capacity efficiency



\[(C_ext/C_int = 0.5). d $\geq$ 5 disfavored by A5 (genericity: extra graviton modes provide no\]

admissibility benefit).



\[The internal sector uses C_int = 12 (dim SU(3)$\times$SU(2)$\times$U(1)), leaving C_ext for geometric degrees\]

\[of freedom. In d = 4, the graviton has 2 propagating DOF, which is the minimum for A4-compatible\]

gravitational records.



\section{Curvature from Capacity Gradients}

Correlation capacity is generally non-uniform. Capacity gradients cause geodesic deviation: paths

converge in low-capacity regions (bottlenecks) and diverge in high-capacity regions. Curvature is

the geometric signature of correlation load---not a postulated feature of spacetime but a derived

consequence of capacity distribution.



The stiffness parameter G (Newton’s constant) sets the scale of geometric response to correlation

load: (curvature change) $\sim$ G $\times$ (load change). G is not a coupling strength but the rate at which

geometry bends under correlation burden. Structurally: $\kappa$ $\sim$ 1/C_* where C_* is the total geometric

capacity (T10).



\section{Einstein’s Equations from $\Delta$_geo Closure}



\subsection{The Geometric Prerequisites}

The February 2026 release closes the complete $\Delta$_geo gap: all five geometric prerequisites

(A9.1--A9.5) are now derived from the axioms, not assumed:



Prerequisite         Content                                       Status



A9.1 Locality        Geometric response depends only on local data

Derived (A1 + FBC)



A9.2 Covariance      Response is coordinate-invariant              Derived (T7B)



A9.3 Conservation Capacity cannot be created or destroyed          Derived (A1)



A9.4 Second-order No higher-derivative instabilities               Derived (A4 + Ostrogradsky)



A9.5 Propagation     Gravitational waves propagate                 Derived (A4: records require DOF)



\subsection{Lovelock Uniqueness}

\textbf{Theorem (Lovelock, 1971) [Import]. In d = 4, the only symmetric, divergence-free tensor}

G_$\mu$$\nu$ that depends on the metric and its first and second derivatives (linearly in second

derivatives) is the Einstein tensor plus cosmological term.



A9.1--A9.5 are precisely Lovelock’s conditions. The unique field equation is:



\[G_$\mu$$\nu$ + Λg_$\mu$$\nu$ = $\kappa$T_$\mu$$\nu$\]



Einstein’s equations are not postulated. They are the unique admissibility-preserving geometric

response in 4D. The equivalence principle follows: because correlation cost depends only on

capacity distribution (not internal properties), all systems follow the same least-cost paths.



Universality of free fall follows from universality of correlation constraints.



\section{The Cosmological Constant as Capacity Residual (T11)}

\textbf{Theorem T11 [P_structural]. Λ is the global capacity residual after all local enforcement}

commitments: Λ $\sim$ (C_total - C_used) / V.



\subsection{Why Λ Is Small}

Almost all correlation capacity has been discharged locally: structure formation, matter, radiation,

horizons, records. Only the asymptotic remainder survives. The remaining scale is set by the

largest admissible interface---the cosmological horizon R_H---giving Λ $\sim$ 1/R_H² $\sim$ H_0². Not

because of vacuum energy, but because R_H is the last interface where correlation enforcement

still matters.



\subsection{Why Λ Is Nonzero}

\[Exact saturation (C_total = C_used) is a measure-zero condition. Any finite universe with finite\]

total capacity generically has a small positive residual. The sign is positive because C_total >

C_used: the universe has not yet fully committed its capacity budget.



\subsection{Why Λ Looks Like Accelerated Expansion}

Local mass-energy is redistributable: it creates curvature gradients and attraction. Global residual

capacity is non-redistributable: it creates uniform pressure and expansion. Λ does not “push space

apart”---it prevents further correlation contraction once global slack is exhausted. Acceleration is

the failure mode of further geometric reorganization.



\section{Dark Matter as Geometric Capacity}

\[The capacity partition C_total = C_int + C_ext identifies a geometric sector C_ext that gravitates (it\]

curves spacetime via T7B) but does not couple to gauge interactions (it has no internal structure).

This is precisely the phenomenology of dark matter: gravitational effects without electromagnetic

or strong-force interactions.



Dark matter is not a particle species. It is the external geometric capacity---the C_ext portion of the

total budget that supports spacetime structure but carries no gauge charges. Its distribution follows

capacity gradients, explaining rotation curves and gravitational lensing without introducing new

particles. Prediction: direct detection experiments should find no dark matter particles.



\section{Where Geometry Fails}

Geometry is valid when correlation cost is locally integrable, capacity gradients are smooth, and

no interface is saturated. Geometry fails at singularities (capacity gradients diverge), horizons

(interfaces persistently saturated), Planck scale (discrete structure dominates smooth

approximation), and topology change (manifold structure inapplicable). These failures are not

pathologies. They are predictions: here is where geometric description exits its domain.



\section{Summary}



\subsection{The Derivation Chain}

\textbf{Result                        Source                                    Status}



Variational dynamics          Smooth enforcement cost in Regime R       [P] (conditional on R)



Euler-Lagrange equations      Extremal paths of cost functional         [P]



Conservation laws             Noether + cost symmetries                 [P]



Correlation distance          Minimal enforcement cost                  [P]



Metric tensor g_$\mu$$\nu$            Non-factorization + polarization (T7B)    [P_structural]



Lorentzian signature          A4 + HKM + Malament                       [P_str] + [I: HKM]



\[d = 4 dimensions              Capacity budget + A4 (T8)                 [P_structural]\]



Einstein equations            A9.1--9.5 + Lovelock (T9_grav)             [P_str] + [I: Lovelock]



Equivalence principle         Universality of correlation cost          [P_structural]



\[$\kappa$ $\sim$ 1/C_*                     Capacity scaling (T10)                    [P_str] (scaling only)\]



Λ $\sim$ 1/R_H²                    Global capacity residual (T11)            [P_structural]



\[Dark matter = C_ext           Geometric capacity partition              [P_structural]\]



\subsection{The Structural Conclusion}

Dynamics and geometry are not fundamental laws but effective representations---the unique

bookkeeping structures available when finite enforceability permits smooth coordination. Together

with Paper 5 (quantum structure), this recovers the full structure of fundamental physics from finite

enforceability---not as postulates, but as the unique bookkeeping available when capacity

constraints permit smooth description. The gravity sector is now structurally closed: metric,

signature, dimension, field equations, cosmological constant, and dark sector all derive from the

axioms.



Appendix A: $\Delta$_geo Closure Summary

The $\Delta$_geo gap encompassed four sub-gaps, all closed in the February 2026 release:



Sub-gap              Content                             Resolution



$\Delta$_ordering           Causal order from A4                A4 irreversibility ⇒ strict partial order. R2/R3 formalized.



$\Delta$_fbc                Finite-bounded cost condition       Lipschitz lemma from A1 + $\\epsilon$-granularity.



$\Delta$_continuum          Smooth manifold structure           Kolmogorov + chartability bridge. Imports Kolmogorov consistency.



$\Delta$_signature          Lorentzian (-,+,+,+)                A4$\to$order + HKM + $\Omega$=1 (Radon-Nikodym uniqueness).



Appendix B: Enforcement Cost vs. Action

Entropy measures irreversibly committed capacity (accounting of commitment). Enforcement cost

measures capacity required for transitions (cost of reallocation). They are related but distinct:

entropy accumulates; enforcement cost is path-dependent. In Regime R, accumulated

\[enforcement cost can be encoded as action: K = $\\int$ L dt. The identification with physical action is\]

exact within the regime and breaks down at saturation. Paper 7 (ACTION) proves that this

accumulated cost has a minimum: the structural origin of Planck’s constant.



Appendix C: Changelog from January 2026 Release

Change                           Section       Details



$\Delta$_geo fully closed               10.1          All A9.1--9.5 derived. No geometric prerequisites remain assumed.



T7B metric derivation            7.1           Non-factorization + polarization identity $\to$ metric. Gap CLOSED.



Lorentzian signature             7.2           A4 $\to$ partial order + HKM/Malament $\to$ (-,+,+,+). CLOSED.



T8 d = 4 derived                 8             d $\leq$ 3 hard-excluded (no graviton DOF). d = 4 minimal. CLOSED.



T11 cosmological constant        11            Λ = global capacity residual. Explains smallness + sign. NEW.



Dark matter section              12            DM = C_ext geometric capacity. No particles predicted. NEW.



Summary table expanded           14.1          Added T10, T11, DM with full epistemic tags



$\Delta$_geo appendix                   App. A        Complete sub-gap resolution table



Date updated                     Header        1/30/2026 $\to$ 2/8/2026




\end{document}
